% !TeX spellcheck = en_US
\documentclass[parskip=full]{report}

\usepackage{amsmath}
\usepackage{amsfonts}
\usepackage{listings}
%\usepackage{beramono}
\usepackage{float}
\usepackage[utf8]{inputenc}
\usepackage[T1]{fontenc}
\usepackage{xcolor}
\usepackage[a4paper, margin={3cm}]{geometry}
\usepackage{hyperref}
\usepackage{graphicx}
\usepackage{svg}
\usepackage{subcaption}
\usepackage{float}
\usepackage{pdfpages}
\usepackage{algorithmicx}
\usepackage[ruled]{algorithm}
\usepackage{algpseudocode}
\usepackage{pgfplots}

\usepackage{tikz}

\usepackage{hyphenat}
\usepackage[english]{babel}
% Carattere monospaziato di default
\renewcommand{\ttdefault}{pcr}

\tikzstyle{block} = [draw, fill=blue!20, rectangle, 
minimum height=3em, minimum width=6em]
\tikzstyle{sum} = [draw, fill=blue!20, circle, node distance=1cm]
\tikzstyle{input} = [coordinate]
\tikzstyle{output} = [coordinate]
\tikzstyle{pinstyle} = [pin edge={to-,thin,black}]

\lstset{
	% wrap long lines on new line
	postbreak=\mbox{\textcolor{red}{$\hookrightarrow$}\space},
	breaklines=true, 
	columns=fullflexible,
	% tab and fonts
	tabsize=2,
	basicstyle=\ttfamily\small,
	% theme
	numbers=left,
	rulecolor=\color{black!30},	
	% UTF8 and escape
	escapeinside={\%TEX}{\^^M},
	inputencoding=utf8,
	extendedchars=true,
	literate={á}{{\'a}}1 {à}{{\`a}}1 {é}{{\'e}}1 {è}{{\`e}}1,
}


% Title Page
\title{
	\includegraphics[width=0.333\textwidth]{assets/unipi1.png} \\
	\textsc{University of Pisa} \\
	\vspace{.5cm}
	Artificial Intelligence and Data Engineering \\
	Cloud Computing \\
	\vspace{2cm}
	{\huge \textsc{K-means} on MapReduce}
}

\author{
	The \textsc{\textbf{Don Matteo}} group: \\
	\vspace{.3cm} \\
	\begin{tabular}{lr}
		Dario Pagani & 585281 \\
		Ricky Marinsalda & 585094 \\
		Giulio Bello & 603078
	\end{tabular}
}

\begin{document}
\maketitle
\tableofcontents


\chapter{Algorithm to select initial centroids}

\section{Idea}

\paragraph{}
To run the \textit{K-means} algorithm is necessary to select $k$ initial centroids, they can either be selected statically by the user --- that is as an algorithm's parameter --- or stochastically from the data; in the latter case they could either be drawn randomly with uniform probability or by employing a probability function that changes the likelihood to draw an element accordingly to certain metrics. For simplicity's sake our implementation chooses the $k$ initial point with equal probability.

\paragraph{}
We used a \textit{Map-Reduce} procedure to draw those $k$ points.

\section{Implementation}

\paragraph{Key}
This \textit{MapReduce} procedure doesn't use a key, so only one reducer will
be spawned by the framework, in our \textit{Hadoop} implementation we used the
\texttt{NullWritable} data-type as output key to reduce traffic, as such object serializes to zero bytes.

\paragraph{Mapper}
The mapper assigns a random label to each data point, this value is not be confused with the \textit{MapReduce}'s key; finally it emits the tuple made of the label and the data point.

\paragraph{Combiner and Reducer}
Then, the combiner sorts the tuples by their label and emits the first $k$
smallest labels and their samples. Finally, the reducer performs the same
operations as the combiner, emitting $k$ values.

\paragraph{Probability}
If the data are split equally among the nodes and the random numbers generator generates all numbers with equal probability, then all dataset's samples have circa equal probability to be drawn.

\paragraph{Complexity}
Let $n = |D|$ be the number of samples, $N$ the number of mapper tasks and $k$ 
the number of samples to draw; then the time complexity is

\[
T \in \mathcal{O} \left(\dfrac{n}{N} \cdot \left(1 + \log(k)\right) + N\cdot k \right)
\]

and the space complexity is

\[
S \in \Theta \left(k\right)
\]

if a sorted data structure is used to store only the first $k$ smallest labels
and their associated data. The I/O complexity is optimal as we perform only linear scans of the file.

\begin{algorithm}[H]
	\caption{Random select}\label{alg:random_map}
	\begin{algorithmic}
		\Require $k \in \mathbb{N}^+$
		\Procedure{Mapper}{nLine, $r$}
			\State $I \gets \Call{rand}{\;}$
				\Comment{Assign a random number to each file's row}
			\State \Return $\left< \texttt{null}, \left< I, r\right> \right>$
				\Comment{Constant key for all lines}
		\EndProcedure
		\vspace{.25cm}
		\Procedure{Combiner}{$S$}
			\State $L$ is a data structure \textbf{ordered} by key $I$
				\Comment{E.g. a binary search tree}
			\State $L \gets \emptyset$
			\While{$S \neq \emptyset$}
				\State $\left< \texttt{null}, \left< I, r\right> \right> \
					\gets \Call{read}{S}$
					\Comment{Read data from the mapper (or combiners)}
				%\State $L \gets L \bigcup \left\{ \left< I, r\right>\right\}$
				\State \Call{insert}{$L$, $\left< I, r\right>$}
				\If{$|L| > k$}
					\Comment{We store at most $k + 1$ elements}
					\State \Call{pop\_last}{L}
					\Comment{We remove the last element (sorted by $I$)}
				\EndIf

				\State \Call{next}{$S$}
			\EndWhile
			\ForAll{$ \left< I, r\right> \in L$}
				\State \Call{emitt}{$\left< \texttt{null}, \left< I, r\right> \right>$}
			\EndFor
		\EndProcedure
		\vspace{.25cm}
		\Procedure{Reducer}{$S$}
			\State \Call{Combiner}{$S$} \Comment{Same as the combiner}
		\EndProcedure
	\end{algorithmic}
\end{algorithm}

\begin{figure}[H]
	\centering
	\caption{Example of generation with $k=4$}
	\label{fig:cluster-example}
	\begin{tikzpicture}
		\begin{axis}[
				width=.8\linewidth,
				title={Example of syntethic clusters ($k=4$), we generate also noise but it was ommited}
			]
			\addplot table [only marks, col sep=comma] {ex0/cluster0.csv};
			\addplot table [only marks, col sep=comma] {ex0/cluster1.csv};
			\addplot table [only marks, col sep=comma] {ex0/cluster2.csv};
			\addplot table [only marks, col sep=comma] {ex0/cluster3.csv};
			
			%\addplot[only marks, mark=*, mark options={fill=green}, mark size=7pt] coordinates {
			%	(7.80555,-19.0115)
			%	(77.8431,77.9533)
			%	(-46.1433,93.5352)
			%	(156.583,156.934)	
			%};
		\end{axis}
	\end{tikzpicture}
\end{figure}

\chapter{K-means}

\section{General idea}

Algorithm \ref{algo:kmeans-reducer} describes how an iteration of K-means is 
performed, 
several iterations are performed until the exit condition is reached. An 
iteration starts with a set of clusters' centroids $C$ and at its end produces 
a new set $C'$, the procedure is said to have converged when:

\begin{equation}
	\forall i \in \left\{1, \cdots, k\right\} \hspace{1cm} |c_i - {c'}_i| < 
	\epsilon
\end{equation}

since --- if bad start centroids were selected --- K-means may take a large 
number of 
iterations to converge, a maximum number of iterations is set as parameter of 
the problem.

\subsection{Average builder}
\label{avgBuild}

Let us define an \textbf{AverageBuilder} as a tuple that contains information 
necessary to compute an average over a set of points ${p_1, \dots, p_s}$:

\begin{align}
	\label{def:avgBuild}
& B = \left\langle \sum_{i = 1}^{s} p_i, s\right\rangle =
\left\langle \Sigma, s\right\rangle
&
p_i \in \mathbb{R}^d \land s \in \mathbb{N}
\end{align}

in this way we can define a function to compute the average:

\begin{equation}
\label{computeAvg}
\text{\textsc{ComputeAverage}}(B) = \dfrac{\Sigma}{s}
= \dfrac{1}{s}\cdot\sum_{i = 1}^{s} p_i
\end{equation}

and another to merge two objects:

\begin{equation}
\label{mergeAvg}
\text{\textsc{MergeBuilder}}(B_a, B_b) =
\left\langle \Sigma_a + \Sigma_b, s_a + s_b\right\rangle
\end{equation}

the latter will be useful to gradually \textit{build} an average from a stream 
of 
data points, like in combiner and reducer.

\subsection{Mapper}
The mapper's role is to assign each point $p \in D$ to the closest centroid 
$c_t \in C$, its output is a key-value pair, where the key represents the 
index of the nearest centroid $t$, and the value is an \emph{AverageBuilder} 
with just one point stored in it.

\paragraph{Input}
The KMeansMapper algorithm takes three parameters as input:

\begin{itemize}
	\item $k$: The key represents the identifier of the data point
	\item $p \in \mathbb{R}^d$: The data point to be assigned to a centroid
	\item $C = \left\{c_1, \dots, c_k\right\}$ : A list of points that are 
	clusters' centroids
\end{itemize}

\paragraph{Output}
The mapper's output is in the form $\left\langle t, \left\langle 
p,1\right\rangle \right\rangle$, where $t$ is the closest centroid's index.

\paragraph{Example Usage}
Suppose we have a K-means clustering problem with three clusters and the following data point and centroids:

\begin{itemize}
	\item	Data point: \textit{p} = $(1, 2, 3)$
	
	\item Centroids: $C =\left\{(4,5,6), (7, 8, 9), (10, 11, 12)\right\} $
\end{itemize}

Applying the KMeansMapper algorithm to the data point and centroids, we get the following result:

\[
\textsc{Mapper}(k, p, C) \longrightarrow
\left\langle 1, \left\langle (1, 2, 3),1\right\rangle \right\rangle
\]

This means that the data point $p$ is closer to the centroid of index 1 $(4, 5, 
6)$ and such as it belong to that centroid's cluster.

\paragraph{Complexity Analysis}
The time complexity of the KMeansMapper algorithm depends on the number of 
centroids in the \textit{centroids}' list, let's denote this number as 
\textit{k}.
A call to the procedure requires a complete visit of the $C$ set, thus 
performing $\Theta(k)$ iterations, each iteration requires the computation of 
the expression $|p - c_i|$ that it's done with $\Theta(d)$ operations; thus 
achieving a total time complexity of

\[
T_\textsc{Mapper} \in \Theta \left(k \cdot d\right)
\]

The space complexity of the algorithm is $\Theta(d)$ since it needs to store an 
additional vector to keep track of the nearest centroid while iterating through 
$C$.

\[
S_\textsc{Mapper} \in \Theta \left(d\right)
\]

\subsection{Combiner and reducer}
The Reducer's role is to compute a new set of clusters $C'$, this is done by spawning a Reducer for each cluster $1, \dots, k$ and each Reducer task computes the centroid's cluster's new central point by averaging all point 
within. In order to reduce the I/O complexity of the MapReduce program, we have 
to introduce a Combiner procedure that reduces the number of element to 
transmit over the network.

As stated before both procedures operate on \textsc{AverageBuilder} objects.

Fundamentally the reducer's and combiner's implementations are almost the same, 
the only difference is that the reducer computes the average of the points by 
performing the operation \textsc{ComputeAverage} as described in Formula 
\ref{computeAvg}.

\paragraph{Input}

The \textbf{KMeansReducer} algorithm takes two parameters as input:
\begin{itemize}
	\item \textit{k}: The cluster's index
	\item \textit{S}: A stream of \textsc{AverageBuilder}s that contains 
	information regarding the cluster's points
\end{itemize}

\paragraph{Output}

The Combiner's output is the same as the mapper's. The Reducer's output is a 
tuple $\left\langle k, {c'}_k \right\rangle$, the second component represents 
the new cluster's centroid for the next iteration of \emph{K-mean}, if the exit 
condition is not met.

%\paragraph{Example Usage}
%
%Suppose we have a K-means clustering problem with three clusters labeled as A, 
%B, and C. After the initial assignment of data points to clusters, we have the 
%following data points assigned to each cluster:

%\begin{itemize}
%	\item Cluster A: $(3, 4, 5)$
%	\item Cluster B: [1, 2]
%	\item Cluster C: [6, 7, 8, 9]
%\end{itemize}

%Applying the \textbf{KMeansReducer} algorithm to each cluster, we get the 
%following new centers:

%\begin{itemize}
%	\item \textbf{reduce}("A", [3, 4, 5]) $\rightarrow$ ("A", 4)
%	\item \textbf{reduce}("B", [1, 2]) $\rightarrow$ ("B", 1.5)
%	\item \textbf{reduce}("C", [6, 7, 8, 9]) $\rightarrow$ ("C", 7.5)
%\end{itemize}
%
%The new centers can then be used in the next iteration of the K-means 
%algorithm to update the cluster assignments based on the proximity to these 
%centers.

\paragraph{Complexity Analysis}

\[
T_\textsc{Combiner} \in \mathcal{O}\left(\dfrac{n}{N} \cdot d\right)
\]

%The time complexity of the \textbf{KMeansReducer} algorithm depends on the 
%size of the $values$ list, which represents the number of data points assigned 
%to a particular cluster. Assuming the number of data points assigned to each 
%cluster is approximately the same, the time complexity can be considered 
%linear, $O(n)$, where $n$ is the number of data points assigned to the cluster.

\[
S_\textsc{Combiner} \in \Theta(d)
\]

The same figure in space complexity applies to the Reducer; while the 
Reducer's time complexity is:

\[
T_\textsc{Reducer} \in \mathcal{O}\left(N \cdot d + d\right) = 
\mathcal{O}\left(N \cdot d\right)
\]

as an instance of Reducer will have to merge the averages coming from $N$ 
Combiner tasks for a certain cluster $c_i$ and as the 
\textsc{ComputeAverage}'s time complexity is of class $\Theta(d)$.

\begin{algorithm}
	\caption{KMeans' iteration for MapReduce}
	\label{algo:kmeans-reducer}
	\begin{algorithmic}[1]
		\Require $B = \left\langle \Sigma, s\right\rangle$ to be an 
		\textsc{AverageBuilder}, as described in Section \ref{avgBuild}
		
		\Require $k \in \mathbb{N}^+$ to be the number of clusters to generate
		
		\Require $C = \{c_1, \dots, c_k\}$ to be the centroids' set
		
		\Procedure{ComputeAverage}{$B$} \Comment{The same as Equation 
			\ref{computeAvg}}
		\State \Return $\dfrac{1}{s}\sum_{i = 1}^{s} p_i$
		\EndProcedure
		
		\vspace{.25cm}
		
		\Procedure{MergeBuilder}{$B_a$, $B_b$} \Comment{The same as Equation
			\ref{mergeAvg}}
		\State \Return $\left\langle \Sigma_a + \Sigma_b, s_a + 
		s_b\right\rangle$
		\EndProcedure
		
		\vspace{.25cm}
		
		\Procedure{Mapper}{offset, $p$, $C$}
		%\State minDistance $\gets +\infty$
		%\State closestCentroid $\gets -1$
		%\ForAll{$c_i \in C$}
		%	\Comment{Assign $p$ to the closest centroid}
		%	\State $d \gets \Call{Distance}{c_i, p}$
		%	\If{$d > $ minDistance}
		%		\State closestCentroid $\gets i$
		%		\State minDistance $\gets d$
		%	\EndIf
		%\EndFor
		
		\State closestCentroid $\gets \underset{c \in C}{\mathrm{argmin}} 
		\left\{\Call{Distance}{c, p}\right\}$ \Comment{$k$ iterations}
		
		\State $B \gets \left\langle p, 1\right\rangle$ \Comment{An 
			AverageBuilder object with just $p$ inside}
		
		\State \Return $\left\langle \text{closestCentroid}, B\right\rangle$
		
		\EndProcedure
		
		\vspace{.25cm}
		
		\Procedure{Combiner}{centroid's index, $S$}
		\State $B \gets \left\langle 0,0 \right\rangle$
		\Comment {Empty average builder}
		
		\ForAll{$B' \in S$}
		\State $B \gets \Call{MergeBuilder}{B, B'}$
		\EndFor
		
		\State \Return $\left\langle k, B\right\rangle$
		\Comment{Returns an \texttt{AverageBuilder}}
		\EndProcedure
		
		\vspace{.25cm}
		
		\Procedure{Reducer}{centroid's index, $S$}
		\State $B \gets \Call{Combiner}{\text{centroid's index},S}$
		\Comment{To merge the builders we re-use the combiner}
		\State \Return \Call{ComputeAverage}{B}
		\Comment{New cluster's centroid}
		\EndProcedure
	\end{algorithmic}
\end{algorithm}

\subsection{General complexity of a single iteration}

\paragraph{}
Hadoop splits the $n$ tuples over multiple tasks on $N$ cluster's nodes; for each tuple, each node has to 
perform a call to the Mapper and each Mapper's result is streamed to the 
Combiner, notice that since the results are streamed to the combiner, it'll 
have to only merge them in time $\Theta(d)$ and since $d \in \mathcal{O}(k\cdot 
d)$ we can omit the term; finally all Hadoop's nodes stream 
their Combiners' results to $k$ Reducers. Those operations have a complexity of:

\begin{equation}
\label{eq:it-t-complexity}
T_\textsc{Kmeans} \in \mathcal{O} \left(
	\dfrac{n}{N}\cdot k \cdot d \; \;+\; \;
	k \cdot \left(N \cdot d\right)
\right)
\simeq
\mathcal{O} \left(
	\dfrac{n}{N}\cdot k \cdot d
\right)
\end{equation}

and 

\[
S_\textsc{Kmeans} \in \Theta \left(k + d\right)
\]

\paragraph{I/O complexity}
In principle, if no combiner is used, the number of bytes $O_\textsc{Kmeans}$ 
to be 
transmitted over the network would be of class $\Theta(n \cdot d)$, since the 
mappers would transmit all of their data tuples to the Reducers; however, if 
the combiner is added to the job, we can reduce greatly this figure, saving a 
lot of I/O over the network:

\[
O_\textsc{Kmeans} \in \mathcal{O} \left(N \cdot k \cdot d\right)
\]

\section{Implementation}

Let's take a look at the main classes and data types we implemented to realize 
our \emph{MapReduce} application on Hadoop.

\paragraph{\texttt{Point} Class}

The \texttt{Point} class represents a point in a $d$-dimensional space:

\[
p = (p_1, \dots, p_d)
\]

It provides the following functionalities:

\begin{itemize}
	\item Computation of the norm
	
	\item Computation of sum between two points and scalar product with a number
	
	\item toString function
	
	\item Implementation of the \texttt{Writable} interface to support 
	serialization and de-serialization
	
\end{itemize}

\subparagraph{I/O}
It takes $(4 + 8d)$ bytes to serialize a \texttt{Point} object.

\paragraph{\texttt{AverageBuilder} Class}
It provides the following functionalities:

\begin{itemize}
	\item Storage of the sum of points and their cardinality (number of points)
	\item Methods to add points to the computation, either individually or by 
	combining with another \texttt{AverageBuilder} object
	\item Computation and retrieval of the average point
	\item Implementation of the \texttt{Writable} interface for serialization 
	and de-serialization
\end{itemize}

The \texttt{Point} class handles point-related operations in the Mapper and 
Reduce classes, while the \texttt{AverageBuilder} class facilitates efficient 
computation of the average of multiple points used in the combiner and reducer.

\subparagraph{I/O}
It takes $(4 + 8d + 8)$ bytes to serialize an \texttt{AverageBuilder} object.

\paragraph{Total I/O}
If no combiner is employed, the application would have to transmit $(4 + 8d + 
8) \cdot n$ bytes over the network, for example a dataset of $n = 4 \times 
10^6$ samples with $d = 6$ features and $k=4$ would generate $240 \text{MB}$ 
\footnote{Not MiB} of traffic; on the other hand, if the combiner is used, it 
would have to 
transmit $(4 + 8d + 8) \cdot N \cdot k$ bytes over the network, using the 
same example and assuming $N=2$, it'd have to send just $480$ bytes, 
shaving off \textbf{two orders of magnitude}!

We can reduce even further those figures by refactoring the implementation a 
little: it's possible to make the $d$ field from \texttt{Point} static --- 
since it's a constant parameter of the algorithm --- saving $4$ bytes; and it's 
possible to use a variable length serializer for the \texttt{cardinality} field 
of \texttt{AverageBuilder} reducing its size to just a couple of bytes.

\chapter{Results}

\section{Synthetic data}

\paragraph{Centroids}
To test our program we had to generate synthetic data. We wrote a program that 
generates $k$ points $c_1, \dots, c_k$, the centroids, by drawing them from a 
random variable 
$C \in \left[-B,B\right]^d$ that has a continuous uniform distribution of 
probability in its domain.

\paragraph{Cluster's data}
For each generated centroid $c_i$, we generate a \emph{Gaussian standard 
vector} $X_i \sim \mathcal{N}_d(\mu, \Sigma)$ with mean $\mu = c_i$ and the 
variance of each component 
$\sigma_{jj}$ with $j = 1, \dots, d$ is 
drawn from a random variable with uniform distribution of probability, this 
allows the program to generate ellipsoidal clusters. Finally, the program 
iterates through all centroids' random vectors $X_0, ..., X_k$ until it 
generates $n$ points, those points are streamed to \texttt{stdout} and written 
to a file.

\paragraph{Noise}
On top of the data we added some noise from a random uniform variable $N \in \left[-B - 	\epsilon, B + \epsilon\right]^d$.

\paragraph{Comparisons}
We had to run the program several times to generate different files with 
different sizes, for comparison's sake it is necessary for the clusters' shapes 
and centroids to be the same every time the program is ran; to achieve this it 
uses the same  \emph{seed} for its pseudo-random number generator used to 
initialize its random variables' parameters.

\paragraph{}
We can see an example of synthetic clusters in Figure \ref{fig:cluster-example}

\section{Benchmark}

\begin{figure}[h]
	\centering
	%\begin{tikzpicture}
	%	\begin{axis}[
	%		xmode=log,
	%		ymode=log,
	%		ymin=300,
	%		ymax=20000,
	%		xlabel={$n$},
	%		ylabel={$t [\text{s}]$}
	%		]
	%		\addplot table [x=size, y=avg, col sep=comma] {../bench/k3mean.csv};
	%		\addplot +[mark=none] coordinates {(3500000, 300) (3500000, 20000)};
	%	\end{axis}
	%\end{tikzpicture}
	\begin{subfigure}[b]{0.425\textwidth}
		\centering
		\begin{tikzpicture}
		\begin{axis}[
				ymin=20,
				ymax=70,
				xmode=log,
				%ymode=log,
				xlabel={$n$},
				ylabel={$t [\text{s}]$},
				legend style={at={(0.05,0.95)},anchor=north west}
			]
			\addplot table [col sep=comma,x=size,y=time,]{../bench/data.3.4.csv}; 
			\addlegendentry{$k = 4$}
			\addplot table [col sep=comma,x=size,y=time,]{../bench/data.3.8.csv};
			\addlegendentry{$k = 8$}
			\addplot table [col sep=comma,x=size,y=time,]{../bench/data.3.16.csv};
			\addlegendentry{$k = 16$}
			
			\addplot +[mark=none, color=red] coordinates {(4500000, 20) (4500000, 70)};
		\end{axis}
		\end{tikzpicture}
		\caption{Execution times with $d = 3$}
	\end{subfigure}
	\hfill
		\begin{subfigure}[b]{0.425\textwidth}
		\centering
		\begin{tikzpicture}
			\begin{axis}[
				ymin=20,
				ymax=70,
				xmode=log,
				%ymode=log,
				xlabel={$n$},
				ylabel={$t [\text{s}]$},
				legend style={at={(0.05,0.95)},anchor=north west}
				]
				\addplot table [col sep=comma,x=size,y=time,]{../bench/data.6.4.csv}; 
				\addlegendentry{$k = 4$}
				\addplot table [col sep=comma,x=size,y=time,]{../bench/data.6.8.csv};
				\addlegendentry{$k = 8$}
				\addplot table [col sep=comma,x=size,y=time,]{../bench/data.6.16.csv};
				\addlegendentry{$k = 16$}
				
				\addplot +[mark=none, color=red] coordinates {(2500000, 20) (2500000, 70)};
			\end{axis}
		\end{tikzpicture}
		\caption{Execution times with $d = 6$}
	\end{subfigure}
	\vskip\baselineskip
	\begin{subfigure}[b]{0.425\textwidth}
		\centering
		\begin{tikzpicture}
			\begin{axis}[
				ymin=20,
				ymax=70,
				xmode=log,
				%ymode=log,
				xlabel={$n$},
				ylabel={$t [\text{s}]$},
				legend style={at={(0.05,0.95)},anchor=north west}
				]
				\addplot table [col sep=comma,x=size,y=time,]{../bench/data.12.4.csv}; 
				\addlegendentry{$k = 4$}
				\addplot table [col sep=comma,x=size,y=time,]{../bench/data.12.8.csv};
				\addlegendentry{$k = 8$}
				\addplot table [col sep=comma,x=size,y=time,]{../bench/data.12.16.csv};
				\addlegendentry{$k = 16$}
				
				
				\addplot +[mark=none, color=red] coordinates {(1500000, 20) (1500000, 70)};
			\end{axis}
		\end{tikzpicture}
		\caption{Execution times with $d = 12$}
	\end{subfigure}
	\hfill
		\begin{subfigure}[b]{0.425\textwidth}
		\centering
		\begin{tikzpicture}
			\begin{axis}[,
				xlabel={$k$},
				ylabel={$t [\text{s}]$},
				legend style={at={(0.05,0.95)},anchor=north west}
				]
				\addplot table [col sep=comma,x=k,y=time,]{../bench/data.vark.3.csv}; 
				\addlegendentry{$d = 3$}
				\addplot table [col sep=comma,x=k,y=time,]{../bench/data.vark.6.csv};
				\addlegendentry{$d = 6$}
				\addplot table [col sep=comma,x=k,y=time,]{../bench/data.vark.12.csv};
				\addlegendentry{$d = 12$}
			\end{axis}
		\end{tikzpicture}
		\caption{Execution times with $n = 4 \times 10^6$}
	\end{subfigure}
	\caption{Execution time of a single iteration of K-means. Right of red line $|F| > 128$ MiB}
\end{figure}

\paragraph{Methodology}
We ran our application on synthetic datasets of various combinations of $n, d, k$, we used a Bash script to log information about each execution's exit code, execution time, number of iterations and total time to run all iterations' jobs; this procedure was done a couple of times for each dataset. Finally, we computed the average time to execute a K-means iteration for each execution and the average for all runs of the same dataset $(n,k,d)$.

\paragraph{Results}
%We observe, as expected, a linear growth of the execution time with respect to 
%$n$; but only after the file's size exceeds the chunk limit of $128\text{MB}$, 
%with $d = 6$ it's when $n > 3 \times 10^6$
%\footnote{
%	Note that we're using 
%	logarithmic scales, so a linear relation will 
%	be represented as a line of slope $1$ and offset $log_{10}(m)$, where $m$ 
%	is the slope of the linear relation.
%}
We can observe that, after a certain $\bar{n}$, all execution times start to show a linear dependency in $n$; a similar observation holds true with $k$, however --- strangely enough --- we observe that when $d$ is increased, the execution time slightly drops with respect to $k$, this might be due to CPUs' pipelining capabilities, since most of the $d$-dependent operations occur in loops of such length.

\paragraph{Execution time}
We can observe that, independently from the dataset's dimension $n$ and dimensionality $d$, it takes a certain time to bootstrap the Hadoop job; in particular we have to read and write $k$ centroids to HDFS. We've observed that such time is $1.25k + 20$ seconds, if we remove such time from our figures, we see that application's complexity (Equation \ref{eq:it-t-complexity}) holds true for $n > 10^6$ with a constant factor of circa $100 \times 10^{-9}$ seconds.

\paragraph{Parallelism}
Notice how there's a drop in time when the file's size passes $128$ MiB --- that's when $n$ passes the red line --- that's because its contents are spanning over two HDFS's chunks, so Hadoop is very likely to spawn two tasks and assign one chunk to each; in spite of this parallelism, we're not able to observe a halving of in execution time.

\paragraph{Considerations}
We see that the MapReduce framework begins to make sense when we work with large amount of data that can be split over multiple computing nodes; otherwise the system's overhead is so large that the execution time remains --- in our experiments --- constant for smaller datasets. Of course, we benchmarked a system made only of three servers and \textbf{very} limited resources, in a real scenario with more powerful machines the execution time may be better.

\paragraph{Improvements}
Finally, we have to state that it might be possible to improve the procedure's speed by replacing the double precision floating point numbers with just single precision ones or even, with appropriate considerations, with integers as the \textit{ALU} is orders of magnitude faster than the \textit{FPU}; however the latter solution is difficult to implement as we'd have to properly scale values and we'd have to compute roots of integral numbers.

\section{Comparison}

\paragraph{Initial centroids}
We used Skitlearn's \emph{k-means++} procedure to select good initial centroids, then we used them in both Skitlearn's k-means and ours.

\paragraph{Results}
We obtained similar results in most cases, an example is shown in Figure \ref{fig:clusters}; however is some cases we obtained different results, for instance, in a dataset in which two cluster were close together, Skitlearn's implementation was able to properly separate them, while ours was not.

\begin{figure}[h]
	\centering
	\begin{tikzpicture}
		\begin{axis}[width=.98\linewidth,]
			\addplot [only marks, color=purple, mark size=0.7pt] table [only marks, col sep=comma] {ex2/test.0.csv};
			\addplot [only marks, color=cyan, mark size=0.7pt] table [only marks, col sep=comma] {ex2/test.1.csv};
			\addplot [only marks, color=teal, mark size=0.7pt] table [only marks, col sep=comma] {ex2/test.2.csv};
			\addplot [only marks, color=red, mark size=0.7pt] table [only marks, col sep=comma] {ex2/test.3.csv};
			\addplot [only marks, color=lime, mark size=0.7pt] table [only marks, col sep=comma] {ex2/test.4.csv};
			\addplot [mark=+, mark options={fill=green}, mark size=1.5pt] table [only marks, col sep=comma] {ex2/test.N.csv};			
			
			
			\addplot[only marks, mark=*, mark options={fill=green}, mark size=4pt] table [only marks, col sep=comma] {ex2/hadoop.csv};
			
			\addplot [only marks, mark=*, mark options={fill=orange}, mark size=4pt] table [only marks, col sep=comma] {ex2/python.csv};
			
			\legend{,,,,,noise,Our implementation,Skitlearn's implementation};
		\end{axis}
	\end{tikzpicture}
	
	\caption{Comparison between two K-means implementations}
	\label{fig:clusters}
\end{figure}

\end{document}          
